% !TeX root = ../papers/nc-p3.tex
\section{反対階数八の曲線分岐の排除}
\label{sec:rank-eight-branch-exclusion}
\begin{theorem}\label{thm:rank-eight-branch-exclusion}
反対階数8の標準正則窓で射影像が全て高々一次元なら、順逆いずれの辺の射影にも
正次元のファイバーはない。純関係の有無を問わない。
\end{theorem}
\begin{proof}[中央支持4の中心の排除]
系\ref{cor:rank-eight-center-types}によって各位置の順中心は高々一つである。
中央支持4の順中心 $q$ が位置1にあれば、位置0に中央支持2の順中心 $p$ がある。
$p$ の後続直線を $L$、中央支持の消去直線を $N$ とすると $q$ の後続直線は $N$ である。
$p$ の二項対称二乗平面には $N$ の消去二空間の全方向に終わる純関係がある。
それらの方向は $q$ の順支持全体でもあるので、$w(q)$ の末尾からの縮約像に
純関係はない。従ってこの $2\times4\times3$ テンソルの二項鉛筆は $L_3$ である
（系\ref{cor:pure-free-branch-rank-nine}の鉛筆の証明）。
そこで辺 $2\to3$ では第一点が $N$ 上または第二点が $q$ の反対三次曲線上となり、
位置2にある順中心は必ず $N$ に属する。

まず $q\in L$ は不可能である。このとき $L$ の一般点の中央支持は4で、
$p$ の第二曲線の次数は $d=3-f_1\le2$。二直線列の排除より $d=2$ である。
定理\ref{thm:branch-rank-one-residual}の円錐曲線の場合の記号を使う。
$W_{x(t)}^\perp/W_p$ が今は二次元なので、同じ評価により
$w(-,x(s),y(t),-)=[s,t]A(t)$ の一次鉛筆 $A$ は一般階数が高々2である。
円錐曲線への制限で支持を失わない点は第三列の一次式が零となる順中心で、
唯一の $q$ である。従って $A$ の全左・右支持は少なくとも3,4でなければならない。
このような一次鉛筆はない。二つの一般値の行空間は、次元和が4に達して全四空間を張るので
各々二次元で直和である。するとその二行列の和の像は両列空間の和となり、
階数は少なくとも3になってしまう。

従って $q\notin L$ で、位置1に他の順中心はないから、$p$ の第二曲線は
完全三次系による三次曲線である。これは $N$ とは異なり末端直線でもない。
よって第三曲線は $q$ の反対三次曲線であり、$f_2=6-3=3$ より第二曲線上に順中心がある。
その点は先の二者択一より $N$ 上にある。$N$ の全点の前点は $q$ 一つだが、
$p$ の列では $L$ 上に前点を持つ。これは $q\notin L$ に反する。逆向きにも中央支持4はない。
\end{proof}

\begin{proof}[二組の中央二次元中心と三次曲線]
従って全ての中心の中央支持は2である。各順中心 $p$ の二直線を $L_p,N_p$ と書く。
次位置に順中心 $q$ があれば $L_q\cap N_p=\varnothing$ である。実際、対応する
消去二空間の共通方向は、続く二辺の終文字・始文字として連続純関係を与える。
従って $q\in L_p$、第二曲線の次数は2である。第三曲線は動き、第二曲線が
$L_q$ と異なるので $N_q$ 上にある。ところが $f_2=3\cdot2-3-1=2$ より
第二曲線上のある順中心から $N_q$ に後続があり、同じ不交性に反する。
従って順中心の位置は隣り合わない。逆側も同じである。

順中心が一位置だけの場合は定理\ref{thm:branch-centers-opposite-phases}の
次数十の接空間の議論がそのまま適用される。その部分は純関係の不在を使わない。
詳しくは、その中心の列で内部の順中心に達する必要があるため、
最初の三次数は $1,3,6$、第四次数は $10+b_2$。第四曲線が上限10を使い切り、
この位置の順中心に会うのはその曲線だけである。唯一の順中心は終端強連結成分なので
命題\ref{prop:terminal-branch-contact}により、そこに至るパラメータは一つで不分岐である。
しかし五行評価の次数十の曲線はこの被約曲線に一致し、順中心で接次元3であるから矛盾する。
よって順中心は相対する位置0,2の各一点であり、逆側も相対する二位置を占める。

位置0の順中心 $p$ からの列の第二曲線 $C$ は有理正規三次である。
次の曲線 $D$ は直線にはなれない。実際 $p$ の中央支持 $M$ に $C(t)$ を制限した比は
非定数である。$w(p,C(t))$ が与える $2\times3$ の $L_2$ 鉛筆はこの比を $[u:v]$ として
核点 $[v^2:-uv:u^2]$ を持つ。従って $D$ を $p$ の逆支持の双対平面へ射影すると
この滑らかな二次曲線をたどる。射影中心で一定になることは $D$ が動くため不可能である。
ゆえに列は少なくとも第四曲線まで続く。位置2の順中心を $a$ とすれば $D$ は
$L_a$ とは異なり、第四曲線は $N_a$ に等しい。この直線の次は一定である。
実際、$R_3$ の $N_a$ への制限は、二項対称二乗の三関係が消えるため高々三次元である。
一般点の逆支持は3なので三つの残余関係の評価は独立である。
従ってこの点での三テンソルをそれら三評価と固定した三係数テンソルの積の和で書けば、
順支持は固定係数の支持の和で一定である。列の一般点ではそれが三次元なので後続は定点となる。
この終中心を $z$（位置1）とする。逆側の隣の位置には逆中心がなく、逆順でも
第二曲線は完全三次系なので $D$ も有理正規三次である。
添字を二つ進めて他の順中心から得る列も同様で、その終中心を $b$（位置3）とすれば
$b$ の前点直線は $N_p$ である。

$z$ の中央支持の消去直線は $L_a$ であることにも注意する。
実際 $z$ の逆支持は $N_a$ の消去二空間 $M_a$ であり、
$R_3\cap(V_3\otimes M_a)$ は三次元である（逆中心の直線への制限の階数3）。
これには $a$ の二項対称二乗平面が入り、それで尽くされる。
ゆえにこの平面を張る $w(z)$ の縮約の位置3支持は $L_a$ の消去二空間である。
\end{proof}

\begin{proof}[二つの接触と二次式の二次ジェット]
$p$ の列を $p,L_p,C,D,N_a,z$ と書く。$C,D$ のパラメータを移動で揃える。
$L_2$ の核点表示によれば $C$ の $N_p$ からの射影の次数を $k$ として
$D$ の一点からの平面射影の次数（引戻し直線束の次数）は $2k$ となる。
三次埋込みを一点から射影すると、この次数は3または2である。従って $k=1$ で、
$C\cap N_p$ の長さは2である。$N_p$ の全点は $b$ のみを後続とするので、
この交叉の台は $D$ が $b$ に至る一つのパラメータ $B$ である。
従って $N_p$ は $C$ にそこで接する。
逆の列に $z$ の中央支持の記述を使うと、$D\cap L_a$ も長さ2で、
$C$ が $a$ に至るパラメータ $F$ で接する。$B\ne F$ である。
実際 $D(F)\in L_a$ の逆支持は3であるのに $D(B)=b$ の逆支持は2である。

$C(t),D(u)$ への $R_2$ の制限は両三次曲線が全空間を張るため単射で、
六本の独立な $(3,3)$ 式を与える。全て対角で消えるので括弧 $[t,u]$ で割ると、
二組の二項二次式のテンソル積の中の六空間 $J$ が得られる。双対の消去空間 $E$ は三次元である。
二項二次の評価ベクトルを $v(t)$ と書く。$a\times L_a$ 上の消滅と $D$ の接触から、
$t=F$ として $u$ で微分すると、括弧を割った全式は $(F,F)$ で消える。
同様に $N_p\times b$ 上の消滅と $C$ の接触から $(B,B)$ でも消える。従って
\[
                v(F)\otimes v(F),\quad v(B)\otimes v(B)\ \in E .
\]

一方、次数列 $1,3,3,1$ より $C$ 上の順側欠損は $f_2=3$ で、
位置2の唯一の順中心 $a$ に集中する。$C$ は埋込みであるからパラメータは $F$ 一つである。
$J$ の六行を第一変数で評価した六行三列の行列は、この点で階数2、
三次小行列式の局所共通零次数は3である。実際元の関係評価はこれに $[t,u]$ の乗算で
二項二次式を三次式に入れる写像を合成したもので、この乗算は各 $t$ で分裂単射である。
一様化変数 $l$ で単因子は $1,1,l^3$。従って法 $l^3$ で原始な核ベクトル
$k(l)=k_0+lk_1+l^2k_2$ があり、先の接触より $k_0$ は $v(F)$ の非零倍である。
$v(t(l))=e_0+le_1+l^2e_2$ と書くと、二次完全系により三係数は独立である。
$E$ は $v(t(l))\otimes k(l)$ の最初の三係数を含み、
\[
 e_0\otimes k_0,\qquad e_1\otimes k_0+e_0\otimes k_1,\qquad
 e_2\otimes k_0+e_1\otimes k_1+e_0\otimes k_2
\]
がその基底となる。この三空間の非零な可分解元の右因子は全て $k_0$ に比例する。
実際、非零な最大番号の $e_j$ における行は必ず $k_0$ の非零倍である。
しかし第二の接触は異なる右因子 $v(B)$ を持つ可分解元を与え、矛盾する。
\end{proof}
